\section{Anomalies} \label{sec7}
The most immediate application of the technique
developed in the previous sections is calculation of
quantum anomalies which are defined as (non-zero) variations of
quantum effective action with respect to symmetry transformations
of the classical theory. In this section we consider two most
important examples of quantum anomalies. These are conformal (scale)
and chiral anomalies. We also discuss briefly the Index Theorem.


There exists also a broader view on the
anomalies which includes any qualitatively new phenomena of a quantum
theory which are absent in its' classical counterpart. An example is the
so-called dimensional reduction anomaly \cite{Frolov:1999an}.
For two operators $D_1$ and $D_2$ the quantity $\det (D_1) \det (D_2)/
\det (D_1D_2)$ is called the multiplicative anomaly 
\cite{Elizalde:1998nd}
since $\det (D_1D_2)\ne \det (D_1)\det(D_2)$  is an ``anomalous''
property of infinite dimensional operators. It is interesting to note
that the heat kernel is also useful for treatment of these non-standard
anomalies. 

To simplify the discussion in this section we work on manifolds without
boundaries.  
%%%%%%%%
\subsection{Conformal anomaly}\label{s7conf}
Conformal invariance is one of the symmetries which are usually broken
by quantisation. This phenomenon (called conformal or trace anomaly)
is known since mid 1970's 
\cite{Capper:1974ic,Deser:1976yx,Dowker:1976tf,Brown:1977pq,Tsao:1977tj}
(see \cite{Birrell:1982ix,Duff:1994wm} for more extensive literature).

The vacuum polarisation induced by quantum effects is described
by the energy-momentum tensor
\begin{equation}
T_{\mu\nu}=\frac 2{\sqrt{g}} \frac{\delta W}{\delta g^{\mu\nu}}
\,,\label{EMT}
\end{equation}
where $W$ is the quantum effective action calculated on the background
with the metric tensor $g_{\mu\nu}$. Consider the
conformal transformation 
\begin{equation}
g_{\mu\nu}\to e^{2\rho(x)}g_{\mu\nu} \label{contra}
\end{equation}
for an infinitesimal value of the parameter:
$g_{\mu\nu}\to (1+2\delta\rho )g_{\mu\nu}$,
$g^{\mu\nu}\to (1-2\delta\rho )g^{\mu\nu}$. According to (\ref{EMT})
variation of the effective action reads:
\begin{equation}
\delta W = \frac 12 \iM T_{\mu\nu}\delta g^{\mu\nu}
=-\iM T_\mu^\mu \delta \rho \,.\label{convarW}
\end{equation}
It is clear from the equation above that the trace of the energy-momentum
tensor measures conformal non-invariance of the theory. If the classical
action $\mathcal{L}$ is conformally invariant, classical energy-momentum
tensor is traceless. However, even in this case conformal invariance is
typically broken by quantum effects. For this reason, quantum $T_\mu^\mu$
is called {\it trace} or {\it conformal} anomaly.

If the classical action $\mathcal{L}$ is conformally {\it invariant},
the fluctuation operator $D$ is conformally {\it covariant}. This means
that $D$ transforms homogeneously $D\to e^{-2\rho}D$ under 
(\ref{contra})\footnote{Strictly speaking this property holds up to a
similarity transformation $D\to e^{\alpha\rho}De^{-\alpha\rho}$ which
does not change the functional determinant.}.

We restrict ourselves to the one-loop level and employ the zeta function
regularization in which the effective action is expressed through the
zeta function of the operator $D$ (see (\ref{Wzren})). Hence, we have to
study the conformal properties of $\zeta (s,D)$.
One can prove that the variation of the zeta function with
respect to variation of the operator $D$ reads\footnote{\label{ft}
To derive this formula rigorously we have to use the method
of \cite{Atiyah:1980jh,Romanov:1979sc}. We first assume that $s$ is 
sufficiently large to keep us away from the singularities, then
use the Mellin transformation (\ref{zetaheat}), perform the variation,
then perform the transformation back, and then continue the result
to $s=0$. This is a perfectly standard procedure which allows us to 
work with variations of positive integer powers of $D$; see 
\cite{Esposito:1997wt} for further detail.}
\begin{equation}
\delta \zeta (s,D)=-s {\rm Tr}((\delta D)D^{-s-1})\,.
\label{varizeta}
\end{equation}   

Next we use that under infinitesimal conformal transformations
the operator $D$ transforms as
\begin{equation}
\delta D=-2(\delta \rho) D \,.\label{infcontraD}
\end{equation}
This equation yields
\begin{equation}
\delta \zeta (s,D) =2s \zeta (s,\delta \rho,D) \,.\label{convarizeta}
\end{equation}
For the operators which we consider in this section the zeta function
is regular at $s=0$. Consequently, the variation of the effective action
(\ref{Wzren}) reads:
\begin{equation}
\delta W=-\zeta (0,\delta\rho,D)=-a_n(\delta\rho,D)\,,
\label{cvWzren}
\end{equation}
where we have used the relation (\ref{anz0}) between the zeta function
at zero argument and the heat kernel coefficient $a_n$. 

We compare (\ref{convarW}) with (\ref{cvWzren}) to see that
\begin{equation}
T_\mu^\mu (x)= a_n(x,D) \,.\label{Tan}
\end{equation}
Note, that $\zeta (0,D)$ (as well as the
pole term in (\ref{Wzpole})) is conformally invariant. 
Consequently, the conformal anomaly is
not divergent, and (\ref{Tan}) does not contain the normalisation
scale $\mu$.

The anomaly (\ref{Tan}) is defined by the same coefficient $a_n$ as the
divergent part of the effective action in the zeta function regularization.
Important difference is that the divergence is given by the integrated
coefficient $a_n(D)$, while the conformal anomaly is defined by the
localised coefficient $a_n(x,D)$. The use of the integrated coefficient in 
(\ref{Tan}) is insufficient to recover total derivatives in the anomaly.  

As an example, consider quantum scalar field $\phi$ 
in two dimensions\footnote{Much work on conformal anomalies
on two-dimensional manifolds was done in the context of string
theory. Here we like to mention the papers 
\cite{Alvarez:1983zi,Friedan:1984xq,Friedan:1985ii,
Callan:1985ia,Alvarez:1987vr}.}
coupled
to the background dilaton $\Phi$:
\begin{equation}
\mathcal{L}=\int d^2x\sqrt{g} e^{-2\Phi} (\partial_\mu\phi)
(\partial_\nu\phi)g^{\mu\nu} \,.\label{nonminact}
\end{equation}
On dimensional and symmetry grounds the inner product may also
contain an arbitrary function $\Psi$ of the dilaton:
\begin{equation}
\langle \phi_1,\phi_2 \rangle =\int d^2x\sqrt{g} e^{-2\Psi}
\phi_1(x)\phi_2(x) \,.\label{dilaprod}
\end{equation} 
Such couplings and inner products appear, for example, after the
spherical reduction of higher dimensional theories to two
dimension (see the review paper \cite{Grumiller:2002nm} for
more details and further references).

The rescaled field $\tilde\phi =e^{-\Psi}\phi$ possesses the standard
dilaton-independent inner product (\ref{inn-prod-scal}). In terms 
of $\tilde\phi$ the action (\ref{nonminact}) reads  
\begin{eqnarray}
 &  & \mathcal{L}
=\int d^{2}x\sqrt{g}\tilde\phi D \tilde\phi \, ,\label{ehnac} \\
 &  & D=-e^{2(\Psi -\Phi )}g^{\mu \nu }(\nabla _{\mu }\nabla _{\nu }+2(\Psi _{,\mu }-\Phi _{,\mu })\partial _{\nu }\nonumber \\
&  & \qquad \qquad \qquad \qquad 
+\Psi _{,\mu \nu }+\Psi _{,\mu }\Psi _{,\nu }-2\Psi _{,\mu }\Phi _{,\nu }),
\label{ehnA} 
\end{eqnarray}
where comma denotes covariant differentiation with $\nabla$:
$\Psi_{,\nu}=\nabla_\nu\Psi$.
We can bring $D$ to the standard form (\ref{laplaceb})
\begin{equation}
D=-(\hat g^{\mu\nu} \hat\nabla_\mu\hat\nabla_\nu +E) \label{hatlalpl}
\end{equation}
by introducing the effective metric 
$\hat g^{\mu\nu}=e^{2(\Psi -\Phi)}g^{\mu\nu}$ and the covariant derivative
\begin{equation}
\label{ehno}
\hat{\nabla }_{\mu }=\partial _{\mu }+\hat{\Gamma }_{\mu }+\hat{\omega }_{\mu }\, ,\qquad \hat{\omega }_{\mu }=\Psi _{,\mu }-\Phi _{,\mu }\, ,
\end{equation}
 where \( \hat{\Gamma } \) is the Christoffel connection for the metric \( \hat{g} \).
Here the potential $E$ reads 
\begin{equation}
\label{ehnEh}
E=\hat{g}^{\mu \nu }(-\Phi _{,\mu }\Phi _{,\nu }+\Phi _{,\mu \nu })\, .
\end{equation}
Now we combine (\ref{Tan}) with (\ref{a2nobou}) and the definitions
given above to obtain:
\begin{equation}
\label{ehnT}
T_{\mu }^{\mu }=\frac{1}{24\pi }(R-6(\nabla \Phi )^{2}+4\nabla ^{2}\Phi +2\nabla ^{2}\Psi )\,. 
\end{equation}
In the case $\Psi =\Phi$ the expression (\ref{ehnT}) was first obtained in
\cite{Mukhanov:1994ax}, and in the general case in 
\cite{Kummer:1997jr}\footnote{Different values for the numerical coefficients
in the conformal anomaly (\ref{ehnT}) were reported in 
\cite{Bousso:1997cg,Nojiri:1997hx,Mikovic:1998xq,Nojiri:1998sr}.
The reason for these discrepancies was clearly stated by Dowker
\cite{Dowker:1998bm} who confirmed the result (\ref{ehnT}).
More extensive literature can be found in \cite{Grumiller:2002nm}.}.

In four dimensions the conformal anomaly for different spins can be read off
from (\ref{a4in4dim}) with the numerical coefficients given 
in Table \ref{a4table}.
%%%%%%%%
\subsection{Chiral anomaly}\label{s7chiral}
Chiral anomaly was discovered in 1969 by Adler, Bell
and Jackiw \cite{Adler:1969gk,Bell:1969ts} and since
that time plays a crucial role in understanding of the
low energy hadron physics. A detailed introduction
to the field and extensive literature can be found
in \cite{Bertlmann:1996,Zinn-Justin:2002vj}.

The spinor action (\ref{spinS}) with the Dirac operator
given in (\ref{dirop}) is invariant under the gauge
transformations
\begin{eqnarray}
&&\delta_\lambda \psi = -\lambda\psi\,,\qquad
\delta_\lambda \bar \psi = \bar\psi \lambda \,,\nonumber \\
&&\delta_\lambda A_\mu^5 =[A_\mu^5,\lambda ]\,,\nonumber \\
&&\delta_\lambda A_\mu = \partial_\mu \lambda +[A_\mu ,\lambda ]
\label{gautr}
\end{eqnarray}
and local (Euclidean) chiral transformations
\begin{eqnarray}
&&\tilde\delta_\varphi \psi =-i\varphi \gamma^5\psi \,,\qquad
\tilde\delta_\varphi \bar \psi =-i\bar\psi \varphi \gamma^5
\,,\nonumber \\
&&\tilde\delta_\varphi A_\mu^5 =\partial_\mu \varphi +
[A_\mu ,\varphi ] \,,\nonumber \\
&&\tilde\delta_\varphi A_\mu = -[A_\mu^5,\varphi ] 
\label{chitr}
\end{eqnarray}
with anti-hermitian local matrix parameters $\lambda$ and
$\varphi$. The Dirac operator transforms as:
\begin{equation}
\delta_\lambda \Dir=[\Dir ,\lambda ]\,,\qquad
\tilde\delta_\varphi \Dir =i\{ \varphi \gamma^5,\Dir \} \,.
\label{transP}
\end{equation}
The Lie algebra structure of the transformations
(\ref{gautr}) and (\ref{chitr}) is encoded in the following
relations:
\begin{eqnarray}
&&\delta_{\lambda_1}\delta_{\lambda_2} -
\delta_{\lambda_2}\delta_{\lambda_1} 
=\delta_{[\lambda_1 , \lambda_2]} \label{gau-gau}\\
&&\delta_{\lambda}\tilde\delta_{\varphi} -
\tilde\delta_{\varphi} \delta_{\lambda}=
\tilde\delta_{[\lambda ,\varphi ]} \label{gau-chi}\\
&&\tilde\delta_{\varphi_1} \tilde\delta_{\varphi_2}-
\tilde\delta_{\varphi_2}\tilde\delta_{\varphi_1}
=-\delta_{[\varphi_1,\varphi_2]}\,, \label{chi-chi}
\end{eqnarray}
where all transformation parameters are taken at the same
space-time point. If the matrices $\lambda (x)$ and
$\varphi (x)$ belong to a finite-dimensional compact
Lie algebra of some Lie group $\mathcal{G}$, the transformations
(\ref{gautr}) and (\ref{chitr}) generate locally
the group $\mathcal{G}\times \mathcal{G}$ with gauge transformations
belonging the diagonal sub-group.

The gauge invariance can (and should) be retained in quantum
theory while the chiral invariance is typically broken by the
quantisation. Let us study these effects in the zeta function
regularization. Define determinant of the Dirac operator
as a square root of the determinant of the associated
Laplace operator:
\begin{equation}
\log \det \Dir =\frac 12 \log \det \Dir^2
=\frac 12 \log \det D
\label{defdetP}
\end{equation}
The effective action $W$ corresponding to the path integral
(\ref{pathspin}) reads
\begin{equation}
W=-\log Z =-\frac 12 \log \det D
=\frac 12 \zeta (0,D)' +\frac 12 \ln (\mu^2) \zeta (0,D)
\,,
\label{Wspin}
\end{equation}
where $\mu$ is a normalisation scale.

By virtue of (\ref{varizeta})
\begin{equation}
\delta_\lambda \zeta (s,D) =-
\left( s {\rm Tr} ( [D,\lambda ] D^{-s-1}) \right)
=\left( {\rm Tr} ( [D^{-s},\lambda ]) \right)
=0\,.
\end{equation}
This proves gauge invariance of the effective action.

For the chiral transformation we have:
\begin{equation}
\mathcal{A}(\varphi ):=\tilde\delta_{\varphi} W=
-2i{\rm Tr} (\gamma^5\varphi D^{-s})|_{s=0}\,.
\label{chian1}
\end{equation}
$\mathcal{A}(\varphi )$ measures non-invariance of the
effective action with respect to the chiral transformations. It 
is called the {\it chiral anomaly}. One can express $\mathcal{A}$
through the heat kernel coefficients:
\begin{equation}
\mathcal{A}(\varphi )=-2i a_n(\gamma^5\varphi , D) \,,
\label{Ahk}
\end{equation}
where $a_n(\gamma^5\varphi , D)$ is defined as in 
(\ref{asymptotex}) but with a matrix-valued smearing
function instead of the scalar one.

The chiral anomaly $\mathcal{A}(\varphi )$ should satisfy
certain consistency conditions following from the
Lie algebra identities (\ref{gau-gau})-(\ref{chi-chi})
and gauge invariance of the effective action:
\begin{eqnarray}
&&\delta_\lambda \mathcal{ A}(\varphi )=
\mathcal{A}([\lambda ,\varphi ] ) \,, \label{WZcc1}\\
&&\tilde\delta_{\varphi_1} \mathcal{A}(\varphi_2 )-
\tilde\delta_{\varphi_2} \mathcal{A}(\varphi_1 )=0 \,.\label{WZcc2}
\end{eqnarray}
The relations (\ref{WZcc1}) and (\ref{WZcc1}) are called the
Wess-Zumino consistency conditions \cite{Wess:1971yu}.
One can check by a direct calculation that the anomaly
defined in this section indeed satisfies these conditions.
For this reason $\mathcal{A}$ is called the {\it consistent}
anomaly (as opposed to the {\it covariant} anomaly which
we do not consider here).

Let us turn now to calculation of the chiral anomaly
$\mathcal{A}$. Since the smearing function $\gamma^5\varphi$
in the heat kernel coefficient (\ref{Ahk}) is matrix-valued
we need more information than we possess at the moment.
To recover the missing terms one can adopt the strategy of
the paper \cite{Branson:1997ze}. Consider 
${\rm Tr}(Q\exp (-tD))$ with arbitrary matrix-valued
function $Q$ and arbitrary Laplace type operator $D$.
There is an asymptotic expansion
\begin{equation}
{\rm Tr}_{L^2} (Q\exp (-tD)) 
\cong\sum_{k\ge0}t^{(k-n)/2}a_k(Q,D)
\label{Qasymptotex}
\end{equation}
where the coefficients $a_k$ are locally computable. This
means that they can be represented as integrals of local
invariants constructed from $Q$ and local invariants
of the operator $D$. These local invariants enter with
some coefficients (universal constants) which are to be
defined. Since $Q$ does not commute with $E$ and $\Omega$
there are more invariants than before. In the particular
case $Q=I_Vf$ (where $I_V$ is the unit matrix and $f$
is a function) we should recover the old result
(\ref{a0nobou}) - (\ref{a6nobou}). For $k=0,2,4$ this
last requirement is strong enough to recover 
$a_k(Q,D)$ completely. Therefore, calculation of the chiral
anomaly in the dimensions $n=2$ and $n=4$ is quite simple:
take the expression (\ref{a2nobou}) or (\ref{a4nobou})
and substitute $2i\gamma^5\varphi$ for $f$ and
(\ref{oEdir}), (\ref{Omdir}) for $\omega_\mu$, $E$
and $\Omega_{\mu\nu}$. Computation of some $\gamma$-matrix
traces is still required, but nevertheless calculations 
are considerably easier than presented in many papers.
In flat space and with
$A^5=0$ the result is particularly simple:
\begin{eqnarray}
&&n=2:\qquad \mathcal{A}(\varphi )= -\frac 1{2\pi}
\int d^2x {\rm tr}(\varphi \epsilon^{\mu\nu}F_{\mu\nu}) 
\label{Am2} \\
&&n=4:\qquad \mathcal{A}(\varphi )= \frac i{16\pi^2}
\int d^4x {\rm tr}(\varphi \epsilon^{\mu\nu\rho\sigma}
F_{\mu\nu}F_{\rho\sigma} )
\label{Am4}
\end{eqnarray}
where ${\rm tr}$ denotes trace over internal (flavour
or colour) indices. We can easily generalise this result
for arbitrary even dimension $n$. First we observe that
the only way to construct a pseudoscalar density of
appropriate dimension (for $A^5=0$) is to contract
$n/2$ tensors $F_{\mu\nu}$ with the Levi-Civita
tensor $\epsilon^{\mu_1\mu_2\dots\mu_n}$. Such tensor
structure can be produced only by a trace of $\gamma^5$ and
the maximal number
($n$) of the gamma matrices. 
The only invariant having the required  form is
\begin{equation}
\int d^n\,x {\rm tr}(QE^{n/2})\,.  
\label{1917}
\end{equation}
This term does not contain derivatives. Under the trace
$Q$ and $E$ commute. Therefore, we can calculate
the coefficient in front of (\ref{1917}) by considering the particular
case $E\sim {I}_V\times {const.}$. For this simple case
the dependence of the heat kernel on $E$ is given simply by $e^{tE}$.
By picking up an appropriate term in the expansion of the
exponential we find that the coefficient in front of (\ref{1917})
in $a_n$ is $(4\pi )^{-n/2} ((n/2)!)^{-1}$. Next we substitute
$E=\frac 12 \gamma^\mu \gamma^\nu F_{\mu\nu}$ 
and take trace over the spinorial indices to obtain:
\begin{equation}
\mathcal{A}(\varphi )=-\frac{2i (-i)^{n(n-1)/2}}{(4\pi)^{n/2}(n/2)!}
\int d^nx {\rm tr}\left(\varphi \epsilon^{\mu_1\mu_2\dots\mu_n}
F_{\mu_1\mu_2}\dots F_{\mu_{n-1}\mu_n} \right)
\label{Amarb}
\end{equation}
For global chiral transformations ($\varphi = const.$) this formula
was obtained in \cite{Wipf:1995dy}.

As soon as we know the anomaly (\ref{Amarb}) for $A_\mu^5=0$, many
more terms (containing $A_\mu^5$) can be restored by using the 
Wess-Zumino condition (\ref{WZcc2}).

Some comments on the Fujikawa approach 
\cite{Fujikawa:1979ay,Fujikawa:1980eg} 
to the chiral anomaly are in order.
Consider the path integral (\ref{pathspin}).
Since the action is invariant under the chiral rotations,
the only source of the chiral anomaly could be non-invariance
of the path integral measure. The Jacobian factor appearing due
to the change of the spinorial variables (\ref{chitr}) 
can be formally represented as
\begin{equation}
J=\det (1-2i\varphi\gamma^5 )\,. \label{Jac1}
\end{equation}
Then, to the first order of $\varphi$,
\begin{equation}
\mathcal{ A}(\varphi )=
\tilde\delta_\varphi W =-\log J^{-1} \simeq -2i{\rm Tr}(\varphi\gamma^5 ) 
\label{chian2}
\end{equation}
where $J^{-1}$ has appeared due to the negative homogeneity
of the fermionic measure. The operator $\varphi\gamma^5$ is not
trace class on the space of square integrable spinors (though
formally its' trace is zero at every point). The right hand side
of (\ref{chian2}) is therefore ill defined. Fujikawa suggested
to replace (\ref{chian2}) by a regularized expression
\begin{equation}
\mathcal{ A}(\varphi )= -2i \lim_{M\to \infty}
{\rm Tr}\left( \varphi\gamma^5 e^{-D/M^2} \right) \,.
\label{Fuji}
\end{equation}
The chiral anomaly (\ref{Fuji}) has now the heat-kernel
form with the identification $t=M^{-2}$. If we suppose that
all positive powers of the regularization parameter $M$ in
the small $t$ (large $M$) asymptotic expansion are somehow
absorbed in the renormalization, we arrive at 
\begin{equation}
\mathcal{ A}(\varphi ) =-2i a_n (\gamma^5\varphi, D)
\label{Fuji2} 
\end{equation}
that is just the expression (\ref{Ahk}) obtained above
in the zeta function regularization.
This method, based on calculations of the regularized
Jacobians, can be applied to conformal anomaly as well
\cite{Fujikawa:1980vr}\footnote{Applications of this method
to topological anomalies can be found in \cite{Fujikawa:2003gi}.
}.

The most essential ingredient of the anomaly calculation
presented in this section is the homogeneous transformation
low (\ref{transP}) for the Dirac operator. This homogeneity
allowed us to restore the power $-s$ in the transformation
law (\ref{varizeta}) for the zeta function and to obtain a
simple expression for the anomaly (\ref{chian1}). It is clear
therefore that as long as the operators transform homogeneously
we shall obtain relatively simple local expressions for
corresponding anomalies (understood as variations of the
effective action with respect to infinitesimal transformation
of the background fields). This suggests to consider extensions
of the chiral group. For example, one of such extensions 
\cite{Novozhilov:1995he}
identifies the group parameters with the diquark fields.

%%%%%%%
\subsection{Remarks on the Index Theorem}\label{s7index}
This report is mainly devoted to local aspects of the heat kernel
expansion. There is however one ``global'' 
application of the heat kernel which cannot be ignored. This is the
Index Theorem. In this section we briefly sketch formal
mathematical aspects of the index construction and its' relation
to the heat kernel. For more details we refer to 
\cite{Eguchi:1980jx,Gilkey:1995,Esposito:1997mw,Bertlmann:1996,
Zinn-Justin:2002vj}. Physical applications of the index theorem
to gravity, gauge theories and strings are so numerous that
we cannot even mention them all.

The index theorem was first formulated by Atiyah and Singer
\cite{Atiyah:1963a}, and the heat kernel approach appeared
later in \cite{Gilkey:1973,AtiyahBP:1973}. Roughly speaking,
their construction is as follows.
Consider two vector bundles $V_1$ and $V_2$ over a manifold $M$.
Let the operator $P$ map $V_1$ to $V_2$, and let $V_{1,2}$
have non-degenerate inner products $(\ ,\ )_{1,2}$ so that one
can define an adjoint $P^\dag$ by the equation
$(\phi_2,P\phi_1)_2=(P^\dag \phi_2,\phi_1)_1$. As an example,
one can keep in mind $V_1$ (respectively, $V_2$) describing
positive (respectively, negative) chirality spinors. In this
example $P$ is a part of the Dirac operator (see eq.\ (\ref{adiagdir})
below). Let us define two operators $D_1=P^\dag P$ and
$D_2=PP^\dag$ acting on (smooth sections of) $V_1$ and
$V_2$, and let us suppose that $D_1$ and $D_2$ are elliptic.
In this case we deal with an {\it elliptic complex}.

Since elliptic operators may have only a finite number of zero
modes (say, $N_1$ and $N_2$ for $D_1$ and $D_2$ respectively),
we may define the index by the following equations:
\begin{eqnarray}
{\mbox{index}}(P)&=& N_1 - N_2 \nonumber \\
&=& \dim \ker D_1 - \dim \ker D_2 \nonumber \\
&=& \dim \ker P - \dim \ker P^\dag \,.\label{defind}
\end{eqnarray}
We also have the intertwining relations
\begin{equation}
D_1P^\dag =P^\dag D_2 \,,\qquad PD_1=D_2P \,,\label{intertwin}
\end{equation}
which tell us that non-zero eigenvalues of $D_1$ and $D_2$
coincide. Consider now the heat kernels,
\begin{equation}
K(t,D_1)-K(t,D_2)=\sum_{(\lambda_1)} e^{-t\lambda_1}
-\sum_{(\lambda_2)} e^{-t\lambda_2} =
N_1 - N_2 ={\mbox{index}}(P) \label{hkindex}
\end{equation}
(where $\lambda_{1,2}$ are eigenvalues of $D_{1,2}$).

To be more specific, let us suppose that $D_1$ and $D_2$ are of
Laplace type and that boundary conditions and singularities (if any)
are such that the power-law expansion (\ref{asymptotex}) exists.
Then, expanding both sides of (\ref{hkindex}) in a power series
of $t$ one obtains:
\begin{eqnarray}
&&a_k(D_1)-a_k(D_2)=0 \qquad \mbox{for}\quad k\ne n,\nonumber\\
&&a_n(D_1)-a_n(D_2)=\mbox{index}(P) \,.\label{anindex}
\end{eqnarray}
The last equation provides a simple way to calculate the
index from known heat kernel expansion. It also allows to understand
a very important property of the index, namely, that it is a
{\it homotopy invariant} under quite general assumptions.
Indeed, suppose that $P$ depends on a parameter $\alpha$ in 
such a way that all geometric quantities (including
the metric, connection, the matrix potential, the boundary
conditions etc.)
corresponding to $D_1$ and $D_2$ are smooth functions of $\alpha$.
Clearly, no ``essential'' deformations, like changing the order
of the operator, or adding higher derivative terms to the boundary conditions, 
or including
new types of the singularities, are allowed. Under these smoothness
assumptions $a_n(D_1)$ and $a_n(D_2)$ are smooth functions of
$\alpha$ as well. Therefore, the index is also a smooth function
of $\alpha$. Since the index is an integer, it can only be a constant.
Hence, ${\mbox{index}}(P)$ is invariant under the deformations
described above.

The relations (\ref{anindex}) are also useful for calculations
of the heat kernel coefficients as they give restrictions on the
universal constants appearing in front of independent invariants
(see, e.g., \cite{Gilkey:1995}).

Let us consider an example of the {\it spin complex}. Let $M$ be
an even dimensional manifold admitting a spin structure (so that
one can define spinor fields on $M$). In a suitable basis
the chirality matrix may be presented in the diagonal form
(at this point it may be useful to consult sec.\ \ref{s3spin}):
\begin{equation}
\gamma_5 =
\left( \begin{array}{cc}
I & 0 \\
0 & -I \end{array} \right)\,.
\label{diag5} \end{equation}
Let $V_1$ correspond the $+1$ eigenvalue of $\gamma_5$ (positive
chirality), and $V_2$ -- to $-1$. The Dirac operator (\ref{dirop})
anti-commutes with $\gamma_5$ and, therefore, may be represented
in the form:
\begin{equation}
\Dir = \left( \begin{array}{cc} 0 & P^\dag \\ P & 0
\end{array} \right) \,. \label{adiagdir}
\end{equation}
We also have:
\begin{equation}
\Dir^2 = \left( \begin{array}{cc} P^\dag P & 0\\ 0 & PP^\dag 
\end{array} \right) \,. \label{dirsq}
\end{equation}
We conclude that $\mbox{index}(P)$, which now measures difference
between dimensions of the null subspaces of $\Dir$ with positive and
negative chirality, is equal to global chiral anomaly:
\begin{equation}
\mbox{index}(P)={\rm Tr} \left( e^{-tP^\dag P} \right)
-{\rm Tr} \left( e^{-tPP^\dag }\right) =
{\rm Tr} \left( \gamma_5 e^{-t\Dir^2}\right)=a_n (\gamma_5,\Dir^2)
 \label{Dindex}
\end{equation}
Therefore, the index can be calculated by the methods of the previous
section. In particular, the index of the Dirac operator in two and
four dimensions on flat manifolds without background axial vector
fields ($A_\mu^5=0$) can be read off from (\ref{Am2}) and (\ref{Am4})
with $\varphi :=i/2$.
We also see that chiral anomaly provides very important
information on topology of the background.

One can associate an index with any elliptic complex which must not
necessarily contain only two vector bundles. In general,
one has a sequence of bundles $V_p$ and a family of operators
$P_p$ which map (smooth sections of) $V_p$ to (smooth sections of)
$V_{p+1}$. One requires that $P_{p}P_{p+1}=0$ and that the corresponding
``Laplacians'' $D_p=P^\dag_p P_p + P_{p-1}P_{p-1}^\dag$ are elliptic.
Important examples include the de Rham complex (for which $V_p$ consists
of $p$-forms on $M$ and $P$ is the exterior derivative), the
Dolbeaux complex which deals with the forms on a complex manifold,
and the signature complex which treats selfdual and anti-selfdual
forms. More details can be found in \cite{Eguchi:1980jx}.
Another very important construction of this type is the
Witten index \cite{Witten:1982df} which constraints the supersymmetry
breaking.